% BHU PYQ -- Manually transcribed & error-corrected
\documentclass[11pt,a4paper]{article}

\usepackage[a4paper,top=2.5cm,bottom=2.5cm,left=2.8cm,right=2.8cm,headheight=14pt]{geometry}
\usepackage{amsmath,amssymb}
\usepackage[T1]{fontenc}
\usepackage[utf8]{inputenc}
\usepackage{lmodern}
\usepackage{microtype}
\usepackage{fancyhdr}
\usepackage{enumitem}
\usepackage{hyperref}
\usepackage{lastpage}
\usepackage{parskip}

\hypersetup{colorlinks=true,urlcolor=black,linkcolor=black,
  pdftitle={BPT-505: Electromagnetic Theory -- BHU PYQ},pdfauthor={Banaras Hindu University}}

\pagestyle{fancy}\fancyhf{}
\renewcommand{\headrulewidth}{0.4pt}\renewcommand{\footrulewidth}{0.4pt}
\fancyhead[L]{\small   \textit{Banaras Hindu University}}
\fancyhead[R]{\small   \textit{BPT-505: Electromagnetic Theory}}
\fancyfoot[C]{\small Page  \thepage\ of \pageref{LastPage}}


\newlist{parts}{enumerate}{2}
\setlist[parts,1]{label=(\alph*),leftmargin=*,itemsep=5pt,topsep=3pt}
\setlist[parts,2]{label=(\roman*),leftmargin=*,itemsep=3pt,topsep=2pt}

\newcommand{\pts}[1]{\hfill   \textit{[#1]}}


\begin{document}


\begin{center}
  {\large\bfseries BANARAS HINDU UNIVERSITY}\\[2pt]
  {\normalsize B.Sc. (Hons.) Semester V Examination, 2016-17}\\[6pt]
  \rule{0.8\linewidth}{0.5pt}\\[6pt]
  {\large\bfseries Paper No. BPT-505: Electromagnetic Theory}\\[4pt]
  {\normalsize Time: 3 Hours\hfill Maximum Marks: 70}
\end{center}
\rule{\linewidth}{0.4pt}
\smallskip



\noindent   \textit{Instructions: Answer all questions. Symbols have their usual meanings.}
\bigskip




\noindent   \textbf{Question 1.} Answer any   \textbf{four} of the following: \pts{4 $ \times$ 5 = 20}

\begin{parts}
  \item Show that the vector field $\vec{v} = yz \hat{i} + zx \hat{j} + xy \hat{k}$ can be written as the gradient of a scalar field $f(x, y, z)$.
  \item If an infinitely long straight wire carries a uniform line charge density $\lambda$, find the electric field at a distance $r$ from this wire.
  \item If the potential due to a continuous charge distribution is written as $V(\vec{r}) = \frac{1}{4\pi\epsilon_0} \int \frac{\rho(\vec{r}')}{|\vec{r} - \vec{r}'|} \, d \tau'$, obtain $
abla^2 V(\vec{r})$.
  \item A parallel polarized electromagnetic wave is incident from air to polystyrene ($\mu \approx \mu_0$, $\epsilon_r = 2.6$) at the Brewster angle. Determine the transmission angle.
  \item Determine the state of polarization for the wave whose electric field is given by:
    \[ \vec{E}(z,t) = E_0 ( \hat{i} \cos(kz - \omega t) - \hat{j} \sin(kz - \omega t) ) \]
  \item Calculate the intrinsic impedance of free space.
\end{parts}

\medskip\rule{\linewidth}{0.2pt}




\noindent   \textbf{Question 2.}

\begin{parts}
  \item What is a linear dielectric medium? How do Maxwell's equations modify in such a linear dielectric medium? \pts{6}
  \item Does the Poynting vector correspond to a dual character in the energy-momentum conservation laws in electrodynamics? Explain in detail. \pts{6}
\end{parts}

\medskip\rule{\linewidth}{0.2pt}




\noindent   \textbf{Question 3.}

\begin{parts}
  \item In a lossless medium for which intrinsic impedance $\eta = 60\pi$, relative permeability $\mu_r = 1$, and magnetic field:
    \[ \vec{H} = -0.1 \cos(\omega t - z) \hat{i} + 0.5 \sin(\omega t - z) \hat{j} \, \text{A/m} \]
    find the relative permittivity $\epsilon_r$ and the angular frequency $\omega$. \pts{8}
  \item What is a phasor? The electric and magnetic fields in free space are given by $\vec{E} = E_0 \cos(10^8 t + \beta z) \hat{i}$ V/m and $\vec{H} = H_0 \cos(10^8 t + \beta z) \hat{j}$ A/m. Find the phasor form of these fields. \pts{6}
\end{parts}

\medskip\rule{\linewidth}{0.2pt}




\noindent   \textbf{Question 4.}

\begin{parts}
  \item How do the laws of electrostatics enforce boundary conditions on the electric field and the corresponding electrostatic potential? \pts{6}
  \item Why does any net charge reside entirely on the surface of a conductor? Explain. \pts{6}
\end{parts}
\hfill   \textbf{OR}

\begin{parts}
  \item Using the general solution of the Laplace equation in spherical polar coordinates with azimuthal symmetry, determine the potential outside the surface of a hollow empty sphere of radius $R$ if the surface potential is maintained at $V_0( \theta) = V_0 \sin^2( \theta/2)$, where $V_0$ is a constant. \pts{8}
  \item For a uniform magnetic field $\vec{B}$, show that $\vec{A} = -\frac{1}{2} (\vec{r}  \times \vec{B})$ is a valid vector potential such that $
abla \cdot \vec{A} = 0$ and $
abla  \times \vec{A} = \vec{B}$. \pts{6}
\end{parts}

\medskip\rule{\linewidth}{0.2pt}




\noindent   \textbf{Question 5.}

\begin{parts}
  \item A standard air-filled rectangular waveguide with dimensions $a = 8.636\, \text{cm}$ and $b = 4.318\, \text{cm}$ is fed by a $4\, \text{GHz}$ carrier wave. Determine if the $TM_{11}$ mode will propagate. If so, calculate the phase and group velocities of this mode. \pts{8}
  \item Show that a rectangular waveguide does not support $TM_{10}$ and $TM_{01}$ modes. \pts{6}
\end{parts}
\hfill   \textbf{OR}

\begin{parts}
  \item Derive the general expressions of the attenuation constant $\alpha$ and phase constant $\beta$ for electromagnetic wave propagation in conducting media. \pts{8}
  \item Prove the relation $R + T = 1$ for normal incidence of electromagnetic waves at the boundary between two media. \pts{6}
\end{parts}

\vfill
\rule{\linewidth}{0.4pt}

\begin{center}\small
  \href{https://bhu-exam.vercel.app/}{Click here to visit our website}\\[4pt]
  \href{https://me-aryan.vercel.app/}{Click here to visit our contributor}\\[4pt]
  \href{https://quashan-me.vercel.app/}{Click here to visit our contributor}
\end{center}

\end{document}
